\section{Alternativas de Solución}

\subsection*{Alternativa 1: Enfoque Locativo y de Gestión Humana}
Prioriza la optimización del espacio físico y la capacitación del personal, manteniendo la infraestructura tecnológica actual (sin cero papel).
\begin{itemize}
    \item \textbf{Actividades incluidas:}
    \begin{itemize}
        \item \textbf{1.1.1 Suficientes estudios de tráfico:} Estudio de demanda, análisis de accesibilidad y reestructuración de áreas.
        \item \textbf{1.2.1 Priorización del proyecto:} Justificar con indicadores e incluir en el plan anual.
        \item \textbf{2.1.1 Suficiente capacitación técnica:} Plan de capacitación por rol y manuales.
    \end{itemize}
    \item \textbf{Actividades excluidas:} 3.1.1 (Biometría y cero papel) y 4.1.1 (Portal web e IA).
\end{itemize}

\subsection*{Alternativa 2: Transformación Digital Interna y Cero Papel}
Busca modernizar el flujo de trabajo en ventanilla y eliminar los cuellos de botella manuales sin incurrir en altísimos costos de infraestructura física o desarrollos con Inteligencia Artificial.
\begin{itemize}
    \item \textbf{Actividades incluidas:}
    \begin{itemize}
        \item \textbf{1.2.1 Priorización del proyecto.}
        \item \textbf{2.1.1 Suficiente capacitación técnica.}
        \item \textbf{3.1.1 Suficiente digitalización en el área del personal:} Automatizar flujo (cero papel), dotación de biometría y firma digital en tiempo real.
    \end{itemize}
    \item \textbf{Actividades excluidas:} 1.1.1 (Obras locativas profundas) y módulos de IA de la 4.1.1.
\end{itemize}

\subsection*{Alternativa 3: Enfoque Tecnológico Integral}
Representa la automatización total orientada al auto-servicio del ciudadano, requiriendo una madurez tecnológica y financiera avanzada.
\begin{itemize}
    \item \textbf{Actividades incluidas:}
    \begin{itemize}
        \item \textbf{1.2.1 Priorización del proyecto.}
        \item \textbf{3.1.1 Suficiente digitalización en el área del personal.}
        \item \textbf{4.1.1 Presencia de un plan estratégico de tecnologías:} Portal web de citas, ciberseguridad robusta y módulos de Inteligencia Artificial (IA) predictiva para la demanda.
    \end{itemize}
    \item \textbf{Actividades excluidas:} 1.1.1 (Estudios locativos, ya que el enfoque es virtualizar la atención).
\end{itemize}

\subsection*{Alternativa 4: Optimización Administrativa y Tercerización}
Enfocada en resolver la falta de personal y software mediante la contratación de servicios externos (Software as a Service - SaaS) en lugar de desarrollos propios o grandes obras.
\begin{itemize}
    \item \textbf{Actividades incluidas:}
    \begin{itemize}
        \item \textbf{2.1.1 Suficiente capacitación técnica.}
        \item \textbf{3.1.1 Digitalización en el área del personal:} Uso de plataformas para agilizar procesos de contratación de personal de apoyo.
        \item \textbf{4.1.1 (Adaptada):} Contratar plataformas web de terceros pagando suscripción mensual, evitando el desarrollo de un plan estratégico interno complejo.
    \end{itemize}
    \item \textbf{Actividades excluidas:} 1.1.1 (Obras locativas) y el desarrollo propio de sistemas de IA (4.1.1 puro).
\end{itemize}

\section{Análisis Preliminar de Factibilidad}

Se discute en el equipo del proyecto cada alternativa considerando los cinco (5) aspectos solicitados para descartar rápidamente las no viables y ahorrar esfuerzos en el análisis detallado.

\begin{table}[H]
\centering
\caption{Matriz de Evaluación Preliminar de Alternativas}
\label{tab:evaluacion}
\renewcommand{\arraystretch}{1.3}
\begin{tabular}{|p{3.2cm}|p{2.2cm}|p{2.7cm}|p{2.2cm}|p{2.4cm}|p{2cm}|}
\hline
\textbf{Alternativa} & \textbf{Viabilidad Técnica} & \textbf{Aceptabilidad Comunidad} & \textbf{Financiera} & \textbf{Capacidad Institucional} & \textbf{Impacto Ambiental} \\ \hline
\textbf{A1: Locativa} & \textcolor{green}{Alta} & \textcolor{orange}{Media} & \textcolor{orange}{Media} & \textcolor{green}{Alta} & \textcolor{red}{Baja} \\ \hline
\textbf{A2: Digital Interna} & \textcolor{green}{Alta} & \textcolor{green}{Alta} & \textcolor{green}{Alta} & \textcolor{green}{Alta} & \textcolor{green}{Alta} \\ \hline
\textbf{A3: IA / Portal Full} & \textcolor{red}{Baja} & \textcolor{green}{Alta} & \textcolor{red}{Baja} & \textcolor{red}{Baja} & \textcolor{green}{Alta} \\ \hline
\textbf{A4: Tercerización} & \textcolor{green}{Alta} & \textcolor{orange}{Media} & \textcolor{red}{Baja} & \textcolor{orange}{Media} & \textcolor{orange}{Media} \\ \hline
\end{tabular}
\end{table}

\subsection*{Justificación detallada por criterio}

\textbf{1. Viabilidad técnica de construirla o implementarla:}
\begin{itemize}
    \item \textbf{A1 y A2:} Alta. Adecuar espacios, implementar biometría, firma digital e integración RUNT usan tecnologías e ingenierías ya maduras y estandarizadas.
    \item \textbf{A3:} Baja. Implementar IA predictiva requiere una arquitectura de Big Data que supera las capacidades técnicas actuales.
    \item \textbf{A4:} Alta. Usar software de terceros (SaaS) es técnicamente rápido de desplegar.
\end{itemize}

\textbf{2. Aceptabilidad de la alternativa por la comunidad:}
\begin{itemize}
    \item \textbf{A1:} Media. Mejora la comodidad física, pero el trámite sigue siendo presencial, lento y propenso a reprocesos.
    \item \textbf{A2 y A3:} Alta. La automatización del flujo, el cero papel y el agendamiento web impactan directamente en la reducción del tiempo de espera y los reprocesos.
    \item \textbf{A4:} Media. La tercerización masiva suele generar resistencia interna (sindicatos) y desconfianza en el manejo de datos de ciudadanos por terceros.
\end{itemize}

\textbf{3. Financiamiento requerido versus disponible:}
\begin{itemize}
    \item \textbf{A1:} Media. Requiere presupuesto de infraestructura para obras civiles.
    \item \textbf{A2:} Alta (Factible). El costo de hardware biométrico y licencias de gestión documental se ajusta al presupuesto de modernización TIC.
    \item \textbf{A3:} Baja (Inviable). Desarrollos a la medida con IA y ciberseguridad avanzada exceden el presupuesto.
    \item \textbf{A4:} Baja (Inviable). Pagar suscripciones recurrentes (OPEX) a perpetuidad ahoga el presupuesto de funcionamiento del Estado a largo plazo.
\end{itemize}

\textbf{4. Capacidad institucional para ejecutar y administrar:}
\begin{itemize}
    \item \textbf{A1 y A2:} Alta. La entidad tiene experiencia gestionando contrataciones de hardware, capacitaciones y adecuaciones operativas.
    \item \textbf{A3:} Baja. Faltan perfiles especializados (científicos de datos, expertos en ciberdefensa).
    \item \textbf{A4:} Media. Contratar servicios en la nube requiere ajustes legales y normativos complejos en la contratación pública.
\end{itemize}

\textbf{5. Impacto Ambiental:}
\begin{itemize}
    \item \textbf{A1:} Baja (Negativo). No elimina el uso intensivo de papel ni los desplazamientos inútiles.
    \item \textbf{A2 y A3:} Alta (Positivo). El flujo automatizado elimina el papel y reduce la huella de carbono al evitar múltiples viajes a la sede.
\end{itemize}

\section{Conclusión y Descarte de Alternativas}

Basado en el análisis preliminar de factibilidad, se descartan las siguientes alternativas:

\begin{itemize}
    \item \textbf{Alternativa 1 descartada:} No soluciona el problema de los procedimientos manuales (Causa 2) y es ambientalmente ineficiente.
    \item \textbf{Alternativa 3 descartada:} Claramente no viable por limitaciones financieras y de capacidad institucional (requerimientos de IA).
    \item \textbf{Alternativa 4 descartada:} Inviable financieramente a largo plazo y genera fricciones operativas institucionales.
\end{itemize}

\textbf{Alternativa Seleccionada para Análisis Detallado: Alternativa 2 (Transformación Digital Interna y Cero Papel).} Es la única que presenta viabilidad \textit{Alta} en los 5 criterios evaluados. Optimiza el proceso a través de la digitalización del personal, biometría y firma digital, requiriendo un presupuesto manejable y generando impacto positivo inmediato en la calidad de vida del ciudadano sin exceder la capacidad institucional.